\section*{Capítulo \ref{ch:b2:corrosion} --- Corrosión y protección}

\begin{solution}{exo:b2:corrosion:1}
$j = \qty{0.020}{A/m^2}$: $de/dt = 0.020 \times 65.38/(2 \times 96\,485 \times 7.13
\times 10^6) = \qty{9.5e-13}{m/s}$, es decir, \qty{0.030}{mm} por año.
\end{solution}

\begin{solution}{exo:b2:corrosion:2}
Pernos de acero sobre cobre: el hierro, el menos noble, es el ánodo, y los pequeños pernos
sobre un gran cátodo de cobre se corroen deprisa. Clavos de cinc en acero: el cinc se corroe
y protege el acero que lo rodea.
\end{solution}

\begin{solution}{exo:b2:corrosion:3}
En el estrecho hueco entre las chapas y bajo la cabeza del perno, donde el agua
se renueva lentamente y su oxígeno se agota pronto: por \omterm{def:b2:corrosion:differential}{aireación diferencial}
estas zonas confinadas se convierten en ánodos, y las superficies abiertas, en cátodos.
\end{solution}

\begin{solution}{exo:b2:corrosion:4}
La lata de hojalata: el estaño, más noble que el hierro, convierte el acero expuesto en el ánodo de un
gran cátodo. En el cubo galvanizado se corroe el cinc en lugar del
acero.
\end{solution}

\begin{solution}{exo:b2:corrosion:5}
La agitación adelgaza la \omterm{def:b2:current-potential-curves:limiting-current}{capa de difusión}: la meseta del oxígeno sube, la curva anódica
del hierro la corta más arriba, de modo que aumentan tanto el potencial como la \omterm{def:b2:corrosion:uniform}{corriente
de corrosión}. Sin oxígeno, la única reacción catódica que queda en agua neutra
es la lenta reducción del agua: el \omterm{def:b2:corrosion:uniform}{potencial de corrosión} baja y la
\omterm{def:b2:corrosion:uniform}{corriente de corrosión} se vuelve muy pequeña.
\end{solution}

\begin{solution}{exo:b2:corrosion:6}
$t = 0.90 \times 5000 \times 2 \times 96\,485/(65.38 \times 0.10) = \qty{1.33e8}{s}$,
unos 4.2~años.
\end{solution}

\begin{solution}{exo:b2:corrosion:7}
Carga por kilogramo, $nF/M$: cinc \qty{820}{A.h/kg}, magnesio \qty{2.21e3}{A.h/kg},
aluminio \qty{2.98e3}{A.h/kg}. Potenciales estándar: \ce{Zn^2+}/\ce{Zn} $-0.76$,
\ce{Mg^2+}/\ce{Mg} $-2.36$, \ce{Al^3+}/\ce{Al} \qty{-1.68}{V}. En un suelo seco y resistivo
la corriente debe atravesar una gran resistencia: el magnesio, mucho más
por debajo del hierro, ofrece la mayor tensión motriz.
\end{solution}

\begin{solution}{exo:b2:corrosion:8}
$U = 2.0 \times 4.0 + 1.5 = \qty{9.5}{V}$; potencia \qty{19}{W}; por año, $19 \times
8766 = \qty{1.7e5}{W.h}$, \qty{1.7e2}{kW.h}.
\end{solution}

\begin{solution}{exo:b2:corrosion:9}
En la línea de flotación el metal está mojado por agua rica en oxígeno, un cátodo; justo
por debajo, el agua está menos aireada y la zona se convierte en el ánodo de la pila
formada con la línea de flotación: el acero se corroe allí, en una franja.
\end{solution}

\begin{solution}{exo:b2:corrosion:10}
En agua aireada el \omterm{def:b2:corrosion:uniform}{potencial de corrosión} del acero inoxidable cae en su
meseta pasiva: una corriente minúscula, y la película se regenera sola. Los iones cloruro rompen
la película en los puntos débiles; el metal desnudo es allí un pequeño ánodo frente a un
enorme cátodo pasivo, de modo que su densidad de corriente es enorme. Dentro de la picadura los
iones metálicos se hidrolizan y acidifican la disolución, el cloruro entra para compensar
la carga, y la película no puede volver a formarse: la picadura se hace más profunda.
\end{solution}

\begin{solution}{exo:b2:corrosion:11}
El oxígeno se reduce sobre los \qty{22}{cm^2}: \omterm{def:b2:current-potential-curves:ie-curve}{corriente catódica} $22 \times 5 =
\qty{110}{\micro A}$, suministrada en su totalidad por la oxidación de los \qty{2}{cm^2} de acero
cercanos a la junta, $j = \qty{55}{\micro A/cm^2}$. La pérdida es $5.5$ veces la del hierro
a \qty{10}{\micro A/cm^2}, unos \qty{0.64}{mm} por año.
\end{solution}

\begin{solution}{exo:b2:corrosion:12}
El hierro es inmune por debajo de la recta \ce{Fe^2+}/\ce{Fe}, $-0.41 + 0.0296\log 10^{-6} =
\qty{-0.59}{V}$: mantenido por debajo de unos \qty{-0.6}{V}, el acero deja de disolverse. Mucho
más abajo, el agua se reduce cada vez más deprisa sobre el acero (la recta \ce{H+}/\ce{H2} está a
\qty{-0.41}{V} a pH 7 y la \omterm{def:b2:current-potential-curves:overpotential}{sobretensión} sobre acero es de unas décimas de voltio):
se desperdicia corriente y el hidrógeno levanta el recubrimiento.
\end{solution}

\begin{solution}{pb:b2:corrosion:1}
\textbf{1.} \ce{Fe -> Fe^2+ + 2e-}; \ce{O2 + 2H2O + 4e- -> 4OH-}.
\textbf{2.} La reducción del oxígeno, lenta y que alcanza su meseta de difusión,
limita la \omterm{def:b2:current-potential-curves:ie-curve}{corriente catódica}; la curva anódica del hierro la corta en esa meseta.
\textbf{3.} $j = \qty{0.10}{A/m^2}$: $0.10 \times 3.156 \times 10^7/(2 \times 96\,485)
\times 55.845 = \qty{913}{g}$ por metro cuadrado y por año.
\textbf{4.} $913/7.87 \times 10^6 = \qty{1.16e-4}{m}$, \qty{0.12}{mm} por año.
\textbf{5.} $4/0.116 = 34$~años.
\textbf{6.} La corrosión no es uniforme: las rendijas, las diferencias de aireación
entre suelos y los defectos del recubrimiento concentran la \omterm{def:b2:current-potential-curves:ie-curve}{corriente anódica} en
áreas pequeñas, donde las picaduras perforan la pared mucho antes.
\textbf{7.} $E^\circ = \Delta_f G^\circ/(nF)$ para cada par: \ce{Fe} $-0.41$, \ce{Zn} $-0.76$,
\ce{Mg} $-2.36$, \ce{Al} \qty{-1.68}{V}.
\textbf{8.} El cinc, el magnesio y el aluminio, todos por debajo del hierro.
\textbf{9.} El aluminio se pasiva: su película de óxido detiene su disolución y
entonces suministra poca corriente, salvo que se alee para impedir la película (lo que funciona
en agua de mar, no en los suelos).
\textbf{10.} La corriente debe atravesar la resistencia del suelo: la tensión
motriz entre el ánodo y el acero protegido es de unos \qty{1.9}{V} para el
magnesio frente a \qty{0.35}{V} para el cinc (según los potenciales estándar).
\textbf{11.} $3.156 \times 10^7 \times 24.305/(2 \times 96\,485 \times 0.50) =
\qty{7.95}{kg}$ por amperio y por año.
\textbf{12.} $\pi \times 0.30 \times 10\,000 = \qty{9.42e3}{m^2}$.
\textbf{13.} \qty{94.2}{m^2}.
\textbf{14.} $0.020 \times 94.2 = \qty{1.88}{A}$.
\textbf{15.} $1.885 \times 20 \times 3.156 \times 10^7 = \qty{1.19e9}{C}$.
\textbf{16.} $1.19 \times 10^9 \times 24.305/(2 \times 96\,485 \times 0.50) =
\qty{3.00e5}{g}$, \qty{300}{kg} (o $1.885 \times 20 \times 7.95$).
\textbf{17.} $300/15 = 20$ ánodos.
\textbf{18.} La resistencia del suelo limita hasta dónde puede llevar un único ánodo
su corriente a lo largo de la tubería: los ánodos repartidos la protegen de forma uniforme.
\textbf{19.} Un rectificador hace circular corriente desde un ánodo inerte enterrado en un lecho
anódico, a través del suelo, hasta la tubería conectada a su polo negativo.
\textbf{20.} $1.885 \times 5.0 + 1.5 = \qty{10.9}{V}$.
\textbf{21.} $10.9 \times 1.885 = \qty{20.6}{W}$; $20.6 \times 8766 = \qty{1.8e5}{W.h}$,
\qty{180}{kW.h} por año.
\textbf{22.} El agua se reduce en la tubería: se forma hidrógeno, que levanta el recubrimiento y
puede fragilizar el acero, y se desperdicia corriente.
\textbf{23.} $E < -0.41 + 0.0296\log(10^{-6}) = \qty{-0.59}{V}$.
\textbf{24.} Veinte años de protección exigen $\boldsymbol{\approx \qty{3.0e2}{kg}}$
de ánodos de magnesio.
\end{solution}
